\chapter{Suggested Problems}

\begin{remarkbox}
These problems are designed to accompany the notes. They are not meant to be a
separate problem book. Their purpose is to force active use of the definitions,
variational principles, conservation laws, and examples developed in the main
chapters.
\end{remarkbox}

\section{Variational Principle}

\begin{enumerate}
    \item Derive the Euler--Lagrange equation for a real scalar field with
    Lagrangian density
    \[
        \mathcal{L}=\frac{1}{2}\partial_\mu\phi\partial^\mu\phi-V(\phi).
    \]
    Carefully identify the boundary term.

    \item Consider the Lagrangian density
    \[
        \mathcal{L}=\frac{1}{2}(\partial_t\phi)^2
        -\frac{c^2}{2}(\partial_x\phi)^2.
    \]
    Derive the wave equation and state what boundary conditions make the
    variational principle well posed on a finite interval.

    \item Show that adding a total derivative
    \[
        \mathcal{L}\longrightarrow \mathcal{L}+\partial_\mu K^\mu
    \]
    does not change the Euler--Lagrange equations, but changes the boundary term
    in the first variation.

    \item Let
    \[
        S[q]=\int_{t_1}^{t_2}dt\left(\frac{1}{2}m\dot q^2-V(q)\right).
    \]
    Derive Newton's equation from the variational principle and compare the
    calculation with the field-theory derivation.

    \item For a theory with two real fields \(\phi\) and \(\chi\), derive the
    Euler--Lagrange equations for
    \[
        \mathcal{L}=\frac{1}{2}\partial_\mu\phi\partial^\mu\phi
        +\frac{1}{2}\partial_\mu\chi\partial^\mu\chi
        -V(\phi,\chi).
    \]
\end{enumerate}

\section{Noether Currents}

\begin{enumerate}
    \item For a massless real scalar field with
    \[
        \mathcal{L}=\frac{1}{2}\partial_\mu\phi\partial^\mu\phi,
    \]
    show that the constant shift symmetry
    \[
        \phi\longrightarrow \phi+a
    \]
    gives the conserved current \(J^\mu=\partial^\mu\phi\).

    \item For a complex scalar field with
    \[
        \mathcal{L}=\partial_\mu\phi^*\partial^\mu\phi-m^2\phi^*\phi,
    \]
    derive the Noether current associated with global \(U(1)\) phase rotations.

    \item Derive the canonical stress-energy tensor for a real scalar field and
    verify that it is conserved on shell when the Lagrangian density has no
    explicit spacetime dependence.

    \item Show explicitly that the conservation equation
    \[
        \partial_\mu J^\mu=0
    \]
    implies conservation of
    \[
        Q(t)=\int d^d x\,J^0(t,\vect{x})
    \]
    when the flux through spatial infinity vanishes.

    \item Explain why a gauge symmetry should not be interpreted in exactly the
    same way as a global symmetry. Use electromagnetism as the example.
\end{enumerate}

\section{Scalar Fields}

\begin{enumerate}
    \item Starting from
    \[
        (\Box+m^2)\phi=0,
    \]
    insert a plane wave ansatz and derive the dispersion relation
    \[
        \omega^2=|\vect{k}|^2+m^2.
    \]

    \item Compute the phase velocity and group velocity of a massive
    Klein--Gordon mode. Discuss the limits \(|\vect{k}|\gg m\) and
    \(|\vect{k}|\ll m\).

    \item For
    \[
        V(\phi)=\frac{1}{2}m^2\phi^2+\frac{\lambda}{4}\phi^4,
    \]
    determine for which signs of \(m^2\) and \(\lambda\) the potential is bounded
    below.

    \item For the double-well potential
    \[
        V(\phi)=\frac{\lambda}{4}(\phi^2-v^2)^2,
    \]
    find the vacua and compute the mass squared of small fluctuations around
    each vacuum.

    \item Consider a static source coupled to a massive scalar field. Explain why
    the static response is short-ranged and identify the length scale.
\end{enumerate}

\section{Electromagnetism}

\begin{enumerate}
    \item Starting from
    \[
        \vect{E}=-\nabla\Phi-\partial_t\vect{A},
        \qquad
        \vect{B}=\nabla\times\vect{A},
    \]
    show that the two homogeneous Maxwell equations are automatically satisfied.

    \item Prove that the gauge transformation
    \[
        \vect{A}\longrightarrow \vect{A}+\nabla\chi,
        \qquad
        \Phi\longrightarrow \Phi-\partial_t\chi
    \]
    leaves \(\vect{E}\) and \(\vect{B}\) unchanged.

    \item Vary the Maxwell action
    \[
        S[A]=\int d^4x\left(-\frac{1}{4}F_{\mu\nu}F^{\mu\nu}-J_\mu A^\mu\right)
    \]
    and derive
    \[
        \partial_\mu F^{\mu\nu}=J^\nu.
    \]

    \item Show that gauge invariance of the source coupling requires
    \[
        \partial_\mu J^\mu=0.
    \]

    \item Derive the wave equation for \(\vect{E}\) and \(\vect{B}\) in vacuum and
    show that electromagnetic waves are transverse.
\end{enumerate}

\section{Yang--Mills Theory}

\begin{enumerate}
    \item For a non-Abelian gauge field
    \[
        A_\mu=A_\mu^aT_a,
    \]
    derive the component expression
    \[
        F_{\mu\nu}^a=\partial_\mu A_\nu^a-\partial_\nu A_\mu^a
        +g f_{bc}{}^aA_\mu^bA_\nu^c.
    \]

    \item Show that if the Lie algebra is Abelian, the Yang--Mills field strength
    reduces to the Maxwell field strength.

    \item Starting from
    \[
        S_{\mathrm{YM}}=-\frac{1}{4}\int d^{d+1}x\,F_{\mu\nu}^aF^{a\mu\nu},
    \]
    derive the source-free Yang--Mills equation
    \[
        D_\mu F^{\mu\nu}=0.
    \]

    \item Verify the Bianchi identity
    \[
        D_{[\lambda}F_{\mu\nu]}=0
    \]
    using the Jacobi identity for covariant derivatives.

    \item Explain in words why non-Abelian gauge fields self-interact while the
    electromagnetic field does not self-interact at the classical Maxwell level.
\end{enumerate}

\section{Spinor Fields}

\begin{enumerate}
    \item Using the Clifford algebra
    \[
        \{\gamma^\mu,\gamma^\nu\}=2\eta^{\mu\nu}\mathbb{I},
    \]
    show that the Dirac equation implies the Klein--Gordon equation for each
    component of \(\psi\).

    \item Derive the Dirac equation by varying
    \[
        \mathcal{L}_{\mathrm{D}}=\bar\psi(i\gamma^\mu\partial_\mu-m)\psi
    \]
    with respect to \(\bar\psi\).

    \item Derive the adjoint Dirac equation by varying with respect to \(\psi\).

    \item Verify directly that
    \[
        j^\mu=\bar\psi\gamma^\mu\psi
    \]
    is conserved on shell.

    \item Show that minimal electromagnetic coupling gives
    \[
        (i\gamma^\mu D_\mu-m)\psi=0.
    \]
    Explain how local \(U(1)\) symmetry determines the form of \(D_\mu\).
\end{enumerate}

\section{Spontaneous Symmetry Breaking}

\begin{enumerate}
    \item Analyze the minima of
    \[
        V(\phi)=\frac{\lambda}{4}(\phi^2-v^2)^2.
    \]
    Explain why the potential is symmetric but an individual vacuum need not be.

    \item Expand \(\phi=v+\eta\) around one vacuum of the double-well potential
    and compute the quadratic part of the Lagrangian for \(\eta\).

    \item For a complex scalar field with potential
    \[
        V(\phi)=\lambda(\phi^*\phi-v^2/2)^2,
    \]
    describe the set of vacua geometrically.

    \item Explain the difference between explicit symmetry breaking and
    spontaneous symmetry breaking.
\end{enumerate}

\section{Solitons}

\begin{enumerate}
    \item For a static real scalar field in one spatial dimension, show that the
    energy functional is
    \[
        E[\phi]=\int dx\left(\frac{1}{2}(\partial_x\phi)^2+V(\phi)\right).
    \]

    \item For the double-well potential, explain why finite-energy configurations
    must approach vacua as \(x\to\pm\infty\).

    \item Define a topological sector for the one-dimensional double-well model
    using the limiting values of \(\phi\) at spatial infinity.

    \item Explain why a configuration connecting \(-v\) at one end to \(+v\) at
    the other end cannot be continuously deformed into the trivial vacuum while
    keeping finite-energy boundary behavior.
\end{enumerate}

\section{Gravity as Field Theory}

\begin{enumerate}
    \item Compare the role of the metric \(g_{\mu\nu}\) in gravity with the role of
    a scalar field \(\phi\) in scalar field theory. What makes gravity special?

    \item Starting from the idea that matter couples to the metric through the
    stress-energy tensor, explain the meaning of
    \[
        T_{\mu\nu}=-\frac{2}{\sqrt{|g|}}\frac{\delta S_{\mathrm{matter}}}{\delta g^{\mu\nu}}.
    \]

    \item In first-order language, define the torsion two-form
    \[
        T^a=de^a+\omega^a{}_b\wedge e^b
    \]
    and curvature two-form
    \[
        R^{ab}=d\omega^{ab}+\omega^a{}_c\wedge\omega^{cb}.
    \]
    Compare these expressions with Yang--Mills curvature.

    \item Explain why diffeomorphism invariance is more subtle than an ordinary
    internal gauge symmetry.
\end{enumerate}

\section{Boundaries and Charges}

\begin{enumerate}
    \item Starting from a conserved current
    \[
        \partial_\mu J^\mu=0,
    \]
    derive the relation between the time derivative of the charge in a spatial
    region and the flux through the boundary.

    \item Give an example where a current is locally conserved but the charge in a
    finite region changes with time.

    \item Explain why boundary terms in the variation of the action matter for
    defining Hamiltonian generators.

    \item For electromagnetism, explain why gauge transformations that do not
    vanish at the boundary may be associated with charges.

    \item Discuss the statement: ``The physical charge is often more robust than
    the local representative of the Noether current.'' Use improvement terms as
    part of the explanation.
\end{enumerate}

\begin{summarybox}
The problems in this appendix are intended to reinforce the central skills of
classical field theory: varying actions, controlling boundary terms, deriving
currents, identifying symmetries, reading physical content from field equations,
and distinguishing local equations from global charges.
\end{summarybox}
